\section*{Capítulo \ref{ch:b2:microbial-metabolism} --- Metabolismo microbiano y biopelículas}

\begin{solution}{exo:b2:microbial-metabolism:1}
Cianobacteria: foto-lito-autótrofa (luz, agua, $\mathrm{CO_2}$).
\emph{Nitrosomonas}: quimio-lito-autótrofa (oxidación del amonio,
$\mathrm{CO_2}$). \emph{E.~coli} sobre glucosa: quimio-organo-heterótrofa.
Bacteria púrpura no del azufre sobre succinato a la luz:
foto-organo-heterótrofa. \emph{Thiobacillus} sobre sulfuro en la
oscuridad: quimio-lito-autótrofa.
\end{solution}

\begin{solution}{exo:b2:microbial-metabolism:2}
Oxígeno (los milímetros superficiales), nitrato (justo debajo, donde el
oxígeno se ha agotado), óxidos de manganeso y de hierro (la transición
del pardo al gris), sulfato (la capa negra sulfurosa, gruesa en los
sedimentos marinos), $\mathrm{CO_2}$ para la \omterm{def:b2:microbial-metabolism:anaerobic}{metanogénesis} (en lo más
profundo, o justo bajo la superficie en el agua dulce pobre en sulfato).
\end{solution}

\begin{solution}{exo:b2:microbial-metabolism:3}
$5\,\mathrm{C_6H_{12}O_6} + 24\,\mathrm{NO_3^-} + 24\,\mathrm{H^+}
\to 30\,\mathrm{CO_2} + 12\,\mathrm{N_2} + 42\,\mathrm{H_2O}$: 4,8
iones nitrato por glucosa (cada nitrato acepta 5 electrones y cada
glucosa cede 24).
\end{solution}

\begin{solution}{exo:b2:microbial-metabolism:4}
Una comunidad microbiana adherida a una superficie e inmersa en una
matriz de polisacáridos, proteínas y ADN producida por ella misma.
Etapas: adhesión, microcolonia y secreción de matriz, maduración con
torres y canales, dispersión. La placa dental y la caries; las
infecciones de catéteres e implantes (y de los pulmones en la fibrosis
quística); la obstrucción y la corrosión de tuberías y cascos --- o,
de forma útil, los filtros percoladores de la depuración de aguas.
\end{solution}

\begin{solution}{exo:b2:microbial-metabolism:5}
Ocho electrones caen de \qty{-0.42}{V} a \qty{-0.24}{V}:
$\Delta G^{\circ\prime} = -8\times 96.5\times 0.18 =
\qty{-139}{kJ/mol}$ de metano (valor tabulado \qty{-131}{kJ/mol}), es
decir, \qty{35}{kJ} por $\mathrm{H_2}$: como máximo 0,7 ATP por
hidrógeno, en la práctica alrededor de medio. La energía por molécula de
sustrato es tan pequeña que los metanógenos crecen despacio (\omterm{def:b2:unicellular-diversity:division}{tiempos de
generación} de horas a días) y convierten casi todo su sustrato en gas
en lugar de en células.
\end{solution}

\begin{solution}{exo:b2:microbial-metabolism:6}
$\mu = 1.2\,S/(5 + S)$: \qty{0.20}{h^{-1}} ($T = \ln 2/\mu =
\qty{3.5}{h}$), \qty{0.60}{h^{-1}} (\qty{1.2}{h}), \qty{1.09}{h^{-1}}
(\qty{0.64}{h}). Noventa por ciento: $S = 9K_s = \qty{45}{mg/L}$.
\end{solution}

\begin{solution}{exo:b2:microbial-metabolism:7}
$S^* = 0.05\times 0.4/(0.8 - 0.4) = \qty{0.05}{g/L}$; $X^* =
0.4\times(5 - 0.05) = \qty{1.98}{g/L}$; producción $DX^* = 0.4\times
1.98 = \qty{0.79}{g/L/h}$; lavado $D_c = 0.8\times 5/5.05 =
\qty{0.79}{h^{-1}}$.
\end{solution}

\begin{solution}{exo:b2:microbial-metabolism:8}
Energía capturada: $0.25\times 275 = \qty{69}{kJ}$, es decir, 1,4 ATP
por amonio; $12/1.4 = 8.7$, unos nueve iones amonio por carbono --- y
varias veces más si se cuenta el coste de fabricar NADH impulsando
electrones cuesta arriba desde el nitrito. Un heterótrofo recibe su
carbono ya reducido y organizado, y treinta ATP por glucosa; una
nitrificante debe comprar a la vez su energía y su poder reductor a un
dador pobre y construir cada carbono a partir del $\mathrm{CO_2}$, de
modo que convierte en células una fracción ínfima del sustrato que
procesa.
\end{solution}

\begin{solution}{exo:b2:microbial-metabolism:9}
$L_p = \sqrt{2\times 1.2\times 10^{-9}\times 0.2/0.5} =
\sqrt{9.6\times 10^{-10}} = \qty{31}{\micro m}$. Duplicar $c_0$
multiplica $L_p$ por $\sqrt 2$: \qty{44}{\micro m}; cuadruplicar $q$
lo reduce a la mitad: \qty{15}{\micro m}.
\end{solution}

\begin{solution}{exo:b2:microbial-metabolism:10}
De arriba abajo: agua con algas y cianobacterias (fotosíntesis
oxigénica); una banda herrumbrosa de oxidantes del hierro donde el
$\mathrm{Fe^{2+}}$ se encuentra con el oxígeno; un velo blanco de
oxidantes incoloras del azufre (\omterm{def:b2:microbial-metabolism:trophic}{quimiolitótrofas}) donde el sulfuro se
encuentra con el oxígeno; una banda púrpura de bacterias púrpuras del
azufre (\omterm{def:b2:microbial-metabolism:anoxygenic}{fotosíntesis anoxigénica} sobre sulfuro); una banda verde de
bacterias verdes del azufre (lo mismo, con más tolerancia al sulfuro y
menos luz); barro negro de fermentadores y reductoras del sulfato.
(a) Las bacterias púrpuras del azufre necesitan a la vez la luz de
arriba y el sulfuro de abajo, de modo que se sitúan en la interfaz donde
se cruzan los dos gradientes. (b) Las reductoras del sulfato producen
$\mathrm{H_2S}$, que precipita el hierro en forma de FeS negro.
(c) Cerrado para la materia: el azufre, el carbono y el nitrógeno
circulan entre formas oxidadas y reducidas sin salir; abierto para la
energía: entra la luz, sale el calor, y sin luz todos los ciclos se
detienen.
\end{solution}

\begin{solution}{exo:b2:microbial-metabolism:11}
Producción $pn$, pérdida $kA$: $A^* = pn/k$. $A_c = 10^{-8}\times
6.02\times 10^{23} = 6.0\times 10^{15}$ moléculas por litro; $n_c =
kA_c/p = 5\times 10^{-4}\times 6.0\times 10^{15}/5 = 6\times 10^{11}$
por litro, $6\times 10^{8}$ por mililitro. El cultivo denso
($10^{9}$) está por encima del umbral; el agua de mar ($10^{5}$), cuatro
órdenes de magnitud por debajo. En una \omterm{def:b2:microbial-metabolism:biofilm}{biopelícula} las células están a
$10^{11}$ por mililitro de matriz, y la matriz frena la salida de la
señal (lo que equivale a reducir $k$), de modo que una microcolonia de
unos pocos miles de células en un nanolitro alcanza el quórum, mientras
que las mismas células dispersas en un litro nunca lo alcanzarían.
\end{solution}

\begin{solution}{exo:b2:microbial-metabolism:12}
Fijación del nitrógeno (solo los \omterm{def:b2:unicellular-diversity:prokaryote}{procariotas} tienen nitrogenasa): si se
detuviera, el nitrógeno combinado de la biosfera se perdería por
\omterm{def:b2:microbial-metabolism:anaerobic}{desnitrificación} y la producción se hundiría en pocas décadas.
Nitrificación y \omterm{def:b2:microbial-metabolism:anaerobic}{desnitrificación} (solo \omterm{def:b2:unicellular-diversity:prokaryote}{procariotas}): si se detuvieran,
el amonio se acumularía, el nitrato desaparecería y el nitrógeno nunca
volvería a la atmósfera. \omterm{def:b2:microbial-metabolism:anaerobic}{Metanogénesis} y \omterm{def:b2:microbial-metabolism:anaerobic}{respiración anaerobia} (solo
\omterm{def:b2:unicellular-diversity:prokaryote}{procariotas}): si se detuvieran, la materia orgánica de los sedimentos
anóxicos, los rúmenes y las marismas se acumularía sin mineralizarse y
el carbono saldría del ciclo; igualmente, la fotosíntesis oxigénica fue
una invención bacteriana, y sin ella no habría oxígeno que respirar.
Los eucariotas son una incorporación tardía y químicamente limitada a un
planeta de \omterm{def:b2:unicellular-diversity:prokaryote}{procariotas}.
\end{solution}

\begin{solution}{pb:b2:microbial-metabolism:1}
\textbf{1.} $-24\times 96.5\times (0.82 + 0.43) = \qty{-2900}{kJ/mol}$.
\textbf{2.} Nitrato: $-24\times 96.5\times 1.17 = \qty{-2700}{kJ/mol}$;
sulfato: $-24\times 96.5\times 0.21 = \qty{-490}{kJ/mol}$.
\textbf{3.} El oxígeno rinde más, de modo que los aerobios desplazan a
las desnitrificantes en la competencia por la materia orgánica mientras
dura el oxígeno, y el oxígeno reprime las nitrato reductasas; la
\omterm{def:b2:microbial-metabolism:anaerobic}{reducción del sulfato} exige anoxia \emph{y} ausencia de nitrato, así que
su olor indica que el tanque está mal aireado y el nitrato agotado --- y
el $\mathrm{H_2S}$ es tóxico y corroe el hormigón.
\textbf{4.} $0.4\times 2900/50 = 23$ ATP; $0.4\times 2700/50 = 22$;
$0.4\times 490/50 = 4$.
\textbf{5.} Unas seis veces más ($23/4$).
\textbf{6.} Dos ATP por glucosa suponen un rendimiento de alrededor de
\num{0.1}, de modo que el \qty{90}{\%} de la energía del sustrato queda
en los productos (ácidos, alcoholes, $\mathrm{H_2}$), que los metanógenos
convierten en metano con una ganancia igual de pequeña; poca energía,
poca biomasa, y el carbono sale en forma de gas.
\textbf{7.} $10^4\times 0.3 = \qty{3000}{kg}$ de DQO al día.
\textbf{8.} Biomasa $0.4\times 3000 = \qty{1200}{kg}$ de DQO; los
\qty{1800}{kg} de DQO restantes se oxidan y necesitan \qty{1800}{kg}
de $\mathrm{O_2}$.
\textbf{9.} $1800/2 = \qty{900}{kWh}$ al día.
\textbf{10.} $10^4\times 0.04 = \qty{400}{kg}$ de nitrógeno;
$4.57\times 400 = \qty{1830}{kg}$ de $\mathrm{O_2}$.
\textbf{11.} En estado estacionario $\mu = 1/\theta \le \mu_{\max}$:
lavado por debajo de $\theta = 1/\mu_{\max} = \qty{2}{d}$.
\textbf{12.} $D = \qty{0.1}{d^{-1}}$: $S^* = 1\times 0.1/(0.5 - 0.1)
= \qty{0.25}{g/m^3}$.
\textbf{13.} $S^* = 0.1/(0.25 - 0.1) = \qty{0.67}{g/m^3}$, y el
límite de lavado sube a \qty{4}{d}; el operador debe alargar la edad
del fango (purgar menos fango): con $\theta = \qty{20}{d}$, $S^* =
0.05/0.2 = \qty{0.25}{g/m^3}$ de nuevo.
\textbf{14.} $2.86\times 400 = \qty{1144}{kg}$ de DQO al día.
\textbf{15.} Amonio del efluente $0.25\times 10^4 = \qty{2.5}{kg/d}$,
el \qty{0.6}{\%} de la carga; el resto, \qty{397}{kg} de nitrógeno,
sale como $\mathrm{N_2}$: el \qty{99}{\%}.
\textbf{16.} Ahorro: \qty{1144}{kg} de $\mathrm{O_2}$. Demanda total:
$(1800 - 1144) + 1830 = \qty{2490}{kg}$ de $\mathrm{O_2}$ al día.
\textbf{17.} $1200\times 0.6\times 0.35 = \qty{252}{m^3}$ de metano
al día.
\textbf{18.} $252\times 36 = \qty{9070}{MJ} = \qty{2520}{kWh}$;
electricidad $0.35\times 2520 = \qty{880}{kWh}$ al día.
\textbf{19.} Aireación $2490/2 = \qty{1245}{kWh}$ al día: el metano
cubre $880/1245 = \qty{71}{\%}$ de ella.
\textbf{20.} $D_e\,c'' = q$, luego $c = c_0 + ax + qx^2/2D_e$; con
$c(L_p) = 0$ y $c'(L_p) = 0$: $a = -qL_p/D_e$ y $c =
(q/2D_e)(L_p - x)^2$, con $L_p^2 = 2D_ec_0/q$.
\textbf{21.} $L_p = \sqrt{2\times 1.5\times 10^{-9}\times 0.2/0.3} =
\sqrt{2\times 10^{-9}} = \qty{45}{\micro m}$.
\textbf{22.} $45/300 = \qty{15}{\%}$ aerobia. Más abajo, las
desnitrificantes reducen a $\mathrm{N_2}$ el nitrato que difunde desde
la capa superficial nitrificante, usando la materia orgánica que
difunde hacia dentro.
\textbf{23.} Dentro de una misma película, el gradiente de oxígeno crea
una capa aerobia (nitrificación) sobre otra anóxica (\omterm{def:b2:microbial-metabolism:anaerobic}{desnitrificación}),
separadas por unas pocas decenas de micrómetros; el tanque aireado es
óxico en todo su volumen, por lo que la \omterm{def:b2:microbial-metabolism:anaerobic}{desnitrificación} necesita un
tanque anóxico aparte.
\textbf{24.} El cloro reacciona con los polímeros de la matriz y se
consume en ellos antes de llegar a las células profundas; las células
profundas están hambrientas, crecen despacio o están latentes y son
mucho menos sensibles; solo muere la capa superficial, y la película
vuelve a crecer desde abajo.
\textbf{25.} Demanda de oxígeno de \qty{2490}{kg} al día; metano,
\qty{252}{m^3} al día; el metano aporta el \qty{71}{\%} de la
electricidad de aireación (\qty{880}{kWh} de \qty{1245}{kWh}).
\end{solution}
